Il·luminem el càtode d'una cèl·lula fotoelèctrica amb un feix de llum verda de 560~nm de longitud d'ona i observem que s'origina un corrent elèctric. Comprovem que el corrent desapareix quan apliquem una tensió de 0,950~V (potencial de frenada).

\begin{apartat}{1}
Calculeu el treball d'extracció (funció de treball) i el llindar de freqüència del metall del càtode.
\end{apartat}

\begin{apartat}{1}
Expliqueu raonadament si es produirà efecte fotoelèctric quan un feix de llum de longitud d'ona més gran que el llindar de longitud d'ona incideixi sobre el metall. I si la freqüència del feix incident és més gran que el llindar de freqüència del metall?
\end{apartat}

\dades{%
  $m_\text{electró} = 9,11\times 10^{-31}\ \mathrm{kg}$.\\
  $q_\text{electró} = -1,60\times 10^{-19}\ \mathrm{C}$.\\
  $1\ \mathrm{eV} = 1,60\times 10^{-19}\ \mathrm{J}$.\\
  Constant de Planck, $h = 6,63\times 10^{-34}\ \mathrm{J\,s}$.\\
  Velocitat de la llum, $c = 3,00\times 10^{8}\ \mathrm{m\,s^{-1}}$.}
